% Answers to Chapter 7, from lectures/L07/appendix/c_solutions.md; the self-test from
% lectures/L07/README.md.

\facitavsnitt{c7}{Kapitel 7}{Vektorer}
\begin{facit}

\svar{1.1} \sv{a} $P = (12 - 4I) \cdot I = 5$ ger $4I^2 - 12I + 5 = 0$.
\sv{b} $I = 0{,}5\,\text{mA}$ eller $I = 2{,}5\,\text{mA}$
\sv{c} Vid $0{,}5\,\text{mA}$ ligger $10\,\text{V}$ över modulen och förlusten i motståndet är
$1\,\text{mW}$. Vid $2{,}5\,\text{mA}$ ligger $2\,\text{V}$ över modulen och förlusten är
$25\,\text{mW}$. Den första driftpunkten är mest rimlig: modulen får en användbar matningsspänning
och förlusterna är små.

\svar{2.1} \sv{a} $(7;\, 8)\,\text{V}$ \sv{b} $(18;\, 22)\,\text{V}$ \sv{c} $(-1;\, -4)\,\text{V}$
\sv{d} $(-22;\, -28)\,\text{V}$

\svar{2.2} \sv{a} $(0;\, 2)\,\text{A}$ \sv{b} $(2;\, 6)\,\text{A}$ \sv{c} $(-14;\, 16)\,\text{A}$ \sv{d} $(18;\, 0)\,\text{A}$

\svar{2.3} \sv{a} $\sqrt{5}\,\text{V}$ \sv{b} $\sqrt{50} = 5\sqrt{2}\,\text{V}$ \sv{c} $\sqrt{10}\,\text{V}$
\sv{d} $\sqrt{125} = 5\sqrt{5}\,\text{V}$

\svar{2.4} \sv{a} $x = 23$, $y = -5$ \sv{b} $x = -8$, $y = 2$

\svar{2.5} a) och d): $(6;\, -9) = 3\mathbf{I}$ (i fas) och $(-12;\, 18) = -6\mathbf{I}$ (i motfas).

\svar{2.6} $a = 1{,}5$, $b = 3{,}5$

\svar{2.7} \sv{a} $\abs{\mathbf{U}} = \sqrt{32} \approx 5{,}66\,\text{V}$, $\abs{\mathbf{I}} = \sqrt{5} \approx 2{,}24\,\text{A}$
\sv{b} $v_U = 45^{\circ}$, $v_I \approx 116{,}6^{\circ}$ \sv{c} Cirka $71{,}6^{\circ}$

\svar{3.1} \sv{a} $5$ \sv{b} $5$ \sv{c} $5$ \sv{d} $8$ \sv{e} $4$

\svar{3.2} \sv{a} $\overrightarrow{AB} = (4;\, 3)$, $\abs{\overrightarrow{AB}} = 5$
\sv{b} $\overrightarrow{AB} = (3;\, -4)$, $\abs{\overrightarrow{AB}} = 5$
\sv{c} $\overrightarrow{BA} = (-4;\, -3)$: samma längd, motsatt riktning.
\sv{d} $\overrightarrow{BA} = -\overrightarrow{AB}$

\svar{3.3} \sv{a} $10\,\text{V}$ \sv{b} $13\,\text{V}$ \sv{c} $\sqrt{2} \approx 1{,}41\,\text{V}$
\sv{d} $3\sqrt{2} \approx 4{,}24\,\text{V}$ \sv{e} $7\,\text{V}$

\svar{3.4} \sv{a} $45^{\circ}$, kvadrant I, ingen korrigering
\sv{b} $135^{\circ}$, kvadrant II, $+180^{\circ}$
\sv{c} $225^{\circ}$ (eller $-135^{\circ}$), kvadrant III, $+180^{\circ}$
\sv{d} $-45^{\circ}$ (eller $315^{\circ}$), kvadrant IV, ingen korrigering
\sv{e} $36{,}9^{\circ}$, kvadrant I, ingen korrigering

\svar{3.5} \sv{a} $(8{,}66;\, 5{,}00)\,\text{V}$ \sv{b} $(0;\, 5)\,\text{V}$ \sv{c} $(-4;\, 0)\,\text{V}$
\sv{d} $(-1{,}41;\, -1{,}41)\,\text{V}$ \sv{e} $\sqrt{75{,}0 + 25} = \sqrt{100} = 10\,\text{V}$

\svar{3.6} \sv{a} $45^{\circ}$ \sv{b} $90^{\circ}$ \sv{c} Cirka $81{,}9^{\circ}$
\sv{d} $180^{\circ}$ (i motfas)

\svar{3.7} \sv{a} $(12;\, 5)\,\text{V}$ \sv{b} $13\,\text{V}$ \sv{c} Cirka $22{,}6^{\circ}$
\sv{d} Spänningarna är vinkelräta och adderas geometriskt, som kateterna i en rätvinklig
triangel: $\sqrt{12^2 + 5^2} = 13\,\text{V}$.

\svar{3.8} \sv{a} $(30;\, 40)\,\text{V}$ \sv{b} $50\,\text{V}$ \sv{c} Cirka $53{,}1^{\circ}$
\sv{d} $(30;\, 0)\,\text{V}$: spolens och kondensatorns spänningar tar ut varandra, så hela
spänningen ligger över motståndet.

\svar{3.9} \sv{a} $(0{,}6;\, 0{,}8)$ \sv{b} $\sqrt{0{,}36 + 0{,}64} = 1$
\sv{c} $(-0{,}6;\, 0{,}8)$ \sv{d} $(9;\, 12)\,\text{A}$ \sv{e} $(-9;\, -12)\,\text{A}$

\svar{3.10} \sv{a} $(0;\, 0)\,\text{V}$
\sv{b} Fasspänningarna tar ut varandra: deras summa är noll.
\sv{c} Cirka $230\,\text{V}$ och cirka $240^{\circ}$ (eller $-120^{\circ}$)
\sv{d} $(345;\, 199)\,\text{V}$, storleken cirka $398\,\text{V}$

\svar{3.11} \sv{a} $(2;\, -3)\,\text{V}$, $(-8;\, 12)\,\text{V}$ och $(-2;\, 3)\,\text{V}$
\sv{b} $(-8;\, 12)\,\text{V}$ och $(-2;\, 3)\,\text{V}$
\sv{c} $t = -6$
\sv{d} Ja: $\mathbf{U} = 1 \cdot \mathbf{U}$.

\svar{3.12} \sv{a} $(4;\, 3)\,\text{A}$ \sv{b} $5\,\text{A}$ \sv{c} Cirka $36{,}9^{\circ}$
\sv{d} $4\,\text{A}$, alltså $\mathbf{I}_C = (0;\, 4)\,\text{A}$

\svar{3.13} \sv{a} $1000\,\text{VA}$ \sv{b} Cirka $36{,}9^{\circ}$ \sv{c} $0{,}8$
\sv{d} $\mathbf{S} = (800;\, 0)$ och $800\,\text{VA}$: den reaktiva effekten är helt kompenserad.

\svar{3.14} \sv{a} Längden ges av Pythagoras sats, inte av summan: $\abs{(3;\, 4)} = 5$.
\sv{b} Talet multipliceras med båda komponenterna: $(6;\, -2)$.
\sv{c} Vinkeln mellan visarna är skillnaden mellan vinklarna, inte summan: $70^{\circ}$.
\sv{d} Vektorn ligger i kvadrant II, så $180^{\circ}$ ska läggas till: $126{,}9^{\circ}$.
\sv{e} En längd är aldrig negativ: $\abs{-2\mathbf{u}} = 2\abs{\mathbf{u}}$.

\svar{4.1} \sv{a} $\abs{\mathbf{u}} = \abs{\mathbf{v}} = 5$, $v_u \approx 53{,}13^{\circ}$,
$v_v \approx 143{,}13^{\circ}$ (kvadrant II), $\mathbf{u} + \mathbf{v} = (-1;\, 7)$ med beloppet
$\sqrt{50} \approx 7{,}071$ och vinkeln $98{,}13^{\circ}$, $\alpha = 90^{\circ}$: vektorerna är
vinkelräta.
\sv{b} Samma värden som i a).
\sv{c} \code{atan(3.0 / -4.0)} ger $-36{,}87^{\circ}$, vinkeln för $(4;\, -3)$ i kvadrant IV:
kvoten visar inte vilken komponent som var negativ. \code{atan2(y, x)} får komponenterna var för
sig och ser deras tecken.

\svar[Testa dig själv]{T} \sv{1} $\abs{\mathbf{u}} = 5$; vinkeln är cirka $-53{,}1^{\circ}$ (eller
$306{,}9^{\circ}$), kvadrant IV.
\sv{2} $\mathbf{w} = (5, -1)$

\end{facit}
